\section{The finite predictor and the profile functionals}\label{app:predictor}
The formula for \(\cerr\) uses the same predictor at two scales. At finite size we evaluate it
on a specific board. Along the filling process, the rook profile and the operator relating the
four families of lines enter the calculation. Here we fix the normalizations that connect the two
descriptions.

\subsection{The predictor on the full board}
Let $\mathcal D$ be the set of lines in the two diagonal families. For square weights $p_z$,
let $m_\ell=\sum_{z\in\ell}p_z$ and $h_\ell=1-m_\ell$. On the polytope with row and column sums
equal to one, $0\le p_z\le1/4$ and $m_\ell\le1$, we maximize
\[
 H_n(p)=\sum_z[-p_z\log p_z+3(1-p_z)\log(1-p_z)]
            -\sum_{\ell\in\mathcal D} h_\ell\log h_\ell .
\]
We use $0\log0=0$. The canonical maximizer $\pi$ is unique and lies in the interior for all
sufficiently large $n$~\cite[\S3 and (P2.7)]{pc}. In the following expressions we evaluate all
weights and vacancies at $\pi$. Let $A$ be the incidence matrix of squares and lines, let
$g_z=\pi_z(1-\pi_z)/(1-4\pi_z)$, and write
\[
 S_n=\operatorname{diag}(0_R,0_C,h_D,h_A)+A^{\mathsf T}\operatorname{diag}(g)A,\qquad
 L_n=-\tfrac12\sum_z[\log(1-4\pi_z)-4\log(1-\pi_z)].
\]
The pseudodeterminant is the product of the nonzero eigenvalues. The null direction gives the last
normalization in the expression
\[
 \log Q_{\rm pred}(n)=H_n(\pi)+L_n-\tfrac12\log\operatorname{pdet} S_n
                      +\tfrac12\log(4\pi n).
\]
In $4\pi n$, the symbol $\pi$ denotes the circle constant. All logarithms are natural.
This definition fixes the finite predictor whose signed error we study. Its own finite-size
correction relative to $\log(n!)+n\Phi+C_{\rm pred}$ is a different quantity.

\subsection{The two profile functionals}\label{sec:profile-data}
Write \(z=(r,q)\in[0,1]^2\). The limiting density \(\nu(z)\) has uniform row and column marginals.
The diagonal coordinates are \(d=r-q\) and \(a=r+q\), and we integrate along their fibers with
respect to \(dr\). If \(m_f(y)\) is the integral of \(\nu\) along the line with coordinate \(y\)
in the family \(f=D,A\), its final vacancy is \(u_f(y)=1-m_f(y)\). For the residual fraction \(x\)
we set
\[
 U_f(y;x)=u_f(y)+x(1-u_f(y)),\qquad U_R=U_C=x,\qquad
 W(z)=\frac{\nu(z)}{U_D(d;x)U_A(a;x)} .
\]
\paragraph{Rook correction.}
We scale \(v_x(z)=U_D(d;x)U_A(a;x)\) to give it marginals of mass one, obtaining
\(p_x(r,q)=s_x(r)v_x(r,q)s_x(q)\). Symmetry gives the same factor on both sides.
Multiplying \(v_x\) by the separable factor \(\nu/(u_Du_A)\) gives the same scaled kernel.
At \(x=1\) the scaled kernel equals one. Let \(B(r,q)=p_x(r,q)-1\) be a kernel on \(L^2[0,1]\).
Operator products are defined by integration. We set
\[
 T=(I-B^2)^{-1}-I,\qquad H=TB,\qquad
 j(r)=(T-B^2)(r,r),\qquad j_0=\int_0^1j(r)\,dr .
\]
Here \(\langle f\rangle=\int_{[0,1]^2} f(r,q)\,dr\,dq\).
Powers inside these averages are pointwise powers of kernels.
\begin{align*}
 F(x)={}&\tfrac12 j_0^2+\operatorname{tr}(T^2)-\tfrac32j_0-\tfrac32\int_0^1 j(r)^2\,dr\\
 &+\langle j(r)(B+H-T)(r,q)j(q)\rangle\\
 &-\tfrac23\langle T^3\rangle+2\langle B H^2\rangle+\tfrac23\langle H^3\rangle .
\end{align*}
In particular, \(B=0\) at \(x=1\), giving \(F(1)=0\).
\paragraph{Line prefactor.}
For \(x>0\), use a midpoint grid with an even number \(m\) of points per side, and let \(J(z)\)
denote the four lines through square \(z\). There are \(m,m,2m-1,2m-1\) lines of types
\(R,C,D,A\). Set \(\rho_l=1/m\) on rows and columns and \(\rho_l=U_l/(xm)\) on diagonals of
both families. At the intersection \(z\) of two distinct lines, set
\[
 K_{lh}=W(z)\prod_{a\in J(z)\setminus\{l,h\}}U_a,\qquad
 v_R=1,\quad v_C=-1,\quad v_D=v_A=0,\qquad \Pi=vv^{\mathsf T}/2.
\]
The other entries of \(K\) are zero. A diagonal and an antidiagonal intersect on the grid only
when their parities agree. This gives rise to the two sheets of each diagonal family in the continuum.
Writing \(\rho\) also for the diagonal matrix, we define
\begin{align*}
 L_0&=K(I+\rho K+\rho\Pi)^{-1}\rho K-3\Pi,\qquad
 \eta_l=\begin{cases}1,&l\in R,C,\\ x/U_l,&l\in D,A,\end{cases}\\
 s_l&=m^{-1}\sum_{z\in l}\nu(z)W(z),\qquad
 \theta_l=3\eta_l s_l-\sum_h K_{lh}^2\rho_h .
\end{align*}
Let \([f]_m=m^{-2}\sum_z\nu(z)f(z)\). The continuum line prefactor is the limit of
\begin{align*}
 \mathcal V_m(x)={}&4[W^2]_m
 -\tfrac32\Big[W\sum_{l,h\in J(z)}(L_0)_{lh}\Big]_m\\
 &-\tfrac32\Big[W\sum_{l\in J(z)}\theta_l\Big]_m
 +\frac{9}{4m}\sum_l\eta_l s_l^2 .
\end{align*}
The sum over \(l,h\) runs over ordered pairs and includes the diagonal. The value at zero is
obtained by continuation from \(x>0\). Appendix~\ref{app:interior} identifies the insertion
terms of \(d(x)\) using this same profile and the same parity convention.

